\section*{Discussion}
\label{sec:discussion}

\paragraph{Key methodological contributions and relation to prior work.}
Continuum remodeling models driven by mechanical stimuli have a long history, yet practical FE implementations remain challenged by mesh sensitivity, parameter non-uniqueness, and numerical robustness of strongly coupled multiphysics updates \cite{Weinans1992_BehaviorAdaptiveBoneRemodeling,Jacobs1995_NumericalInstabilitiesBone,Meslier2023_UsingFinite,MartinezReina2016_UniquenessBRM}.
Within this landscape, the present work advances a smoothly regularized three-field formulation in which density adaptation is coupled to an evolving orthotropic fabric state and to an explicit nonlocal mechanotransductive stimulus field (\S\ref{sec:contribution-novelty}--\S\ref{sec:strong-forms-nd}).
Compared with classical mechanostat formulations that operate on a local stimulus index, representing mechanotransduction as a diffusion--reaction field connects naturally to nonlocal stimulus concepts proposed in the mechanobiology literature \cite{Giorgio2019_MechanicallyDrivenBiological,Addessi2024_BoneRemodelingApproach}.
At the constitutive level we retain a fabric-guided orthotropic elasticity structure consistent with established fabric--elasticity relationships \cite{Cowin1985_RelationshipBetweenElasticity,Turner1990_FabricDependenceOrthotropic,Zysset1995_AlternativeModelAnisotropic}, but embed it into a log-fabric evolution and evaluation pipeline that is designed to remain well-defined under eigenvalue degeneracy and under bounded anisotropy.
\rev{In that sense the model deliberately stays at the apparent scale, between bottom-up osteon mechanics \cite{Gizzi2022_OsteonMechanics} and homogenized constrained-mixture evolution laws \cite{Vasta2023_ConstrainedMixtureFabric}: we retain a tractable anatomical-mesh closure while acknowledging that lower-scale turnover and damage processes are collapsed into phenomenological constitutive and fabric updates.}

\paragraph{Robust spectral treatment and implications for stability.}
Fabric-guided orthotropy requires repeated spectral evaluation (eigenvalues, eigenprojectors, and spectral functions) both for constitutive evaluation and for constructing the trace-free log-fabric target.
Standard eigendecompositions are well known to become numerically fragile near isotropy or transverse isotropy, where eigenvectors are non-unique and small perturbations can trigger basis flips; in a coupled setting this manifests as non-smooth stresses and discontinuous drivers.
The robust Sylvester-projector construction (\S\ref{sec:projector-robustness}) addresses this by avoiding eigenvectors, enforcing the partition of unity, and smoothly blending to the isotropic and TI limits, yielding bounded and invariant projectors in the degenerate regimes (\cref{fig:sylvester_projectors}).
\rev{The dedicated near-TI constitutive benchmark in \cref{fig:sylvester_stress_benchmark} makes this concrete: for vanishing perturbations around a TI fabric state, naive eigendecomposition-based or only partially regularized projectors retain an $\mathcal{O}(10^{-1})$ relative stress error, whereas the fully blended construction collapses to the correct TI limit at $\mathcal{O}(10^{-8})$. In this sense the robust treatment suppresses artificial stress jumps at the constitutive level that would otherwise contaminate the coupled fixed-point residual and jeopardize convergence of the staggered solver.}
Together with the smoothly regularized clamping of eigenvalue ratios, it yields a constitutive map that is continuous through eigenvalue coalescence and avoids overflow in $\exp(\mathbf{L})$, which is essential when exploring stiff parameter regimes and when aiming for differentiable solvers and calibration workflows.

\paragraph{Role of nonlocal stimulus and regularization: mesh sensitivity and identifiability.}
Mesh-scale ``checkerboarding'' and localization are a recurring failure mode of local remodeling rules, and have motivated node-based discretizations or explicit nonlocal/gradient regularization in earlier studies \cite{Jacobs1995_NumericalInstabilitiesBone,Peerlings1998}.
Here, diffusion in the stimulus, density, and log-fabric submodels provides a controlled length scale that suppresses oscillations without imposing ad hoc post-processing (\S\ref{sec:parameter-governance}, \S\ref{sec:regularization-analysis}).
Although all three appear as diffusion terms, their intent differs: stimulus diffusion $D_S$ serves as a proxy for nonlocal mechanotransduction (signal reach), whereas diffusion in density and log-fabric is used primarily as a phenomenological gradient regularization to mitigate mesh dependence and stabilize the coupled update.
The diffusion sweep on the stiff box benchmark shows that density diffusion $D_\rho$ is the dominant lever for smoothing $\rho$: increasing $D_\rho$ from $0.1D_{\mathrm{ref}}$ to $10D_{\mathrm{ref}}$ reduces the median roughness ratio from $\approx 0.49$ to $\approx 0.04$ and lowers the median coupling subiterations per step from $\approx 15$ to $\approx 12$ (\cref{fig:diffusion_diagnostic}), consistent with elliptic regularization improving contractivity of the coupled map.
Stimulus diffusion $D_S$ provides a weaker but consistent additional smoothing through the mechanostat drive (roughness $\approx 0.31\to 0.18$ over the same factor range) with essentially unchanged coupling cost, whereas varying fabric diffusion $D_A$ has little impact on density roughness, as expected for a regularizer acting primarily on the directional field.
At the same time, introducing nonlocality raises identifiability questions: the mechanotransduction length implied by $D_S$ is inseparable from load magnitudes, kinetics, and time-step size, while $D_\rho$ interacts with kinetics and $\Delta t$ through the implicit density update.
Moreover, treating $(D_S,D_\rho,D_A)$ as independent can be over-parameterized when field-specific calibration data are unavailable.
In terms of density-field smoothness in this benchmark, a practical hierarchy is $D_\rho$ (primary) $>$ $D_S$ (secondary) $>$ $D_A$ (tertiary); importantly, this ranking reflects the chosen roughness metric and does not diminish the mechanobiological role of $D_S$ in setting a signal length scale.
Accordingly, in applications without dedicated calibration it is reasonable to prescribe $D_S$ (and $D_A$) via a single target length $\ell=\alpha h$ (\S\ref{sec:parameter-governance}) and then use the smallest $D_\rho$ that suppresses mesh-scale artifacts.
The mesh-based prescription $\ell=\alpha h$ provides a pragmatic route to mesh-insensitive predictions, but relating $\ell$ to microstructural or cellular length scales remains an important open problem for future experimental calibration and uncertainty quantification.

\paragraph{Solver strategy: convergence behavior, scalability, and failure modes.}
The coupled update is advanced by a partitioned block Gauss--Seidel strategy stabilized by Anderson (Pulay/DIIS) acceleration with safeguards (\S\ref{sec:fp-aa}).
The discretization study indicates consistent behavior of the coupled state under refinement: spatial $L^2$ pair errors exhibit slopes between first and second order (depending on the field), and temporal refinement yields approximately first-order behavior (\cref{fig:individual_convergence}), matching expectations for CG1 spatial discretization combined with Backward--Euler time integration.
The fixed-point solver tests highlight a key practical regime split.
For the well-conditioned \textbf{PhysioSetup}, both damped Picard and Anderson-accelerated iterations converge below $\mathrm{tol}=10^{-6}$ (final $r_{\mathrm{Pic}}\approx 5\times 10^{-7}$) with comparable cost (33.6~s vs.\ 44.8~s); for the challenging \textbf{StiffSetup}, pure Picard stagnates at a large residual ($r_{\mathrm{Pic}}\approx 4.7\times 10^{-1}$) and wastes work before termination (2.1~min), whereas Anderson mixing restores convergence to tolerance in the same configuration (50.7~s; \cref{fig:anderson_comparison}).
The diagnostic view (\cref{fig:anderson_diagnostic}) and parameter sweep (\cref{fig:anderson_sweep}) illustrate the expected trade-offs: larger history sizes can reduce iterations but deteriorate Gram conditioning and trigger restarts, while regularization stabilizes the least-squares solve at the cost of weaker acceleration \cite{Pulay1980_ConvergenceAcceleration,Walker2011_AndersonAcceleration,Toth2015_ConvergenceAnalysis}.
From a failure-mode perspective, the stiff regime emphasizes that robustness hinges on (i) event-driven restarts and step limiting, (ii) sufficient regularization in the model (diffusion and bounded kinetics), and (iii) reasonable time steps; when these conditions are violated the coupled map can lose contractivity and the iteration can stall or hit the subiteration cap.
The modular partitioned architecture remains attractive for large-scale settings because each block can be solved efficiently with MPI-parallel iterative linear solvers and preconditioned independently; however, further gains will likely require block-specific multigrid strategies and/or Jacobian-free Newton--Krylov alternatives when pushing to more extreme parameter regimes or tighter tolerances.

\paragraph{Interpretation of femur case study and sensitivity to loading assumptions.}
The proximal femur results support two complementary observations.
First, a compartment-aware reference stimulus is not merely a modelling embellishment: sweeping the cortical--trabecular ratio $r=\psi_{\mathrm{ref}}^{\mathrm{cort}}/\psi_{\mathrm{ref}}^{\mathrm{trab}}$ improves similarity to the population-template CT reference across global and compartment-restricted metrics (\cref{fig:psi_ref_ratio_analysis}).
Over the tested range $r\in[0.33,3]$, the total RMSE decreases from 0.61 to 0.56~g/cm$^3$ (cortical RMSE: 0.61 to 0.55~g/cm$^3$).
\rev{To put these RMSE values in perspective, the inter-subject variability of the mapped CT densities in the HFValid cohort, characterized by the $\pm 2\sigma$ band around the population mean, spans approximately 0.3--0.5~g/cm$^3$ in the trabecular core and up to 0.8~g/cm$^3$ in the cortical shell. The observed total RMSE of 0.55--0.61~g/cm$^3$ is thus of a magnitude comparable to the natural inter-subject spread, indicating that the model predictions fall within the population envelope, although the error remains large relative to the trabecular-only variability.}
While the best-RMSE case in this discrete sweep lies at the upper end of the tested range (here $r\approx 3$; \cref{fig:psi_ref_ratio_analysis}), \rev{this value is specific to the current population template, load ensemble, and parameter set and should not be interpreted as a universal optimum; in particular, the discrete coarse sweep does not resolve the sensitivity surface finely enough to claim a sharp optimum.}
The more general qualitative point is that a substantial cortical--trabecular set-point contrast ($r>1$) is needed in this continuum setting to maintain a thin cortical shell under the same loads.
\rev{This interpretation must also be separated from the image-to-model representation of the cortical shell itself, because CT-based cortical mapping can materially alter proximal-femur strain predictions even before remodeling is considered \cite{Falcinelli2020_CorticalAssignment}.}
Because the reference is a mapped population template (not subject-specific longitudinal ground truth), these values quantify proximity to that reference rather than clinical error.
This aligns with prior evidence that mechanostat thresholds and ``set points'' are not universal constants and may differ between compartments or cell populations \cite{Frost1992_SetPointsOsteoporosis,Mullender1998_SetPointTrabecularArchitecture,Yang2021_CancellousGreaterThreshold,Schulte2013_LocalMechanicalStimuli}.
\rev{The additional $\delta_0$--$\kappa_S$ sweep (\cref{fig:mechanostat_sensitivity}) sharpens this point: in the explored neighborhood of the femur default, the saturation width $\kappa_S$ is the more influential local calibration lever, while $\delta_0$ has only a weak first-order effect on the final density field.}
At the same time, the final density distributions reveal a strong accumulation near the imposed bounds, suggesting that the bounded-kinetics closure and the rescaled CT target together can yield bimodal spectra; addressing this may require refined remodeling laws (e.g.\ turnover- or surface-based kinetics) and/or data-driven calibration of bounds and reference densities.
Second, the hip-force direction is an important uncertainty source: a modest frontal-plane offset changes not only medial--lateral redistribution but also global density trajectories, with measurable effects on similarity-to-template metrics and pattern correlation (\cref{fig:alpha_front_summary}).
In the explored range $\Delta\alpha_{\mathrm{front}}\in[-15^\circ,25^\circ]$, the total RMSE decreases from 0.582 to 0.531~g/cm$^3$, so the lowest RMSE occurs at the upper sweep boundary ($\Delta\alpha_{\mathrm{front}}=25^\circ$). Within the physiologic $\pm 13^\circ$ band discussed in \S\ref{sec:alpha-front-sweep}, however, the RMSE varies only from 0.563 to 0.581~g/cm$^3$ and the Pearson correlation from 0.602 to 0.654, while the medial--lateral center-of-density shifts by about 1.26~mm across the full sweep.
Because the present load model uses a small set of representative cases with reduced muscle sets, these sensitivities should be interpreted as a lower bound on the uncertainty introduced by subject- and activity-dependent loading; conversely, the sweep demonstrates that incorporating load-direction uncertainty is necessary for credible case-study comparisons and for robust parameter identification in future calibration workflows.

\paragraph{Limitations and outlook.}
Several limitations delimit the current scope.
Biologically, the model operates at the apparent-density continuum level and does not represent explicit trabecular architecture, surface remodeling dynamics, damage, or biochemical regulation; the diffusion length scales and kinetic gains remain phenomenological calibration parameters.
Methodologically, the accelerated large-$\Delta t$ benchmarks are intentionally used as stress tests for nonlinear contractivity and do not resolve stimulus transients, while the femur case study is referenced to a population-template CT field and simplified loading rather than subject-specific longitudinal data.
\rev{Recent advances in CT-based cortical property assignment \cite{Falcinelli2020_CorticalAssignment} and refined assessment of femoral mechanical competence from clinical CT \cite{Schileo2024_CTFemurAssessment} suggest that tighter integration of patient-specific CT pipelines with the present remodeling framework is a promising direction, in particular for replacing the population-template reference with subject-specific baselines.}
Numerically, the partitioned approach avoids monolithic Jacobians but does not provide the unconditional robustness of a fully coupled Newton method, and extreme parameter combinations can still lead to stall or slow convergence.
\rev{The new near-TI constitutive benchmark isolates the degeneracy mechanism at the material-point level, but a full FE-level ablation of robust versus naive projector handling inside the complete coupled box benchmark remains future work.}
Future work should therefore focus on (i) calibration and uncertainty quantification of nonlocal lengths and mechanostat parameters under realistic load ensembles, (ii) tighter integration with musculoskeletal or instrumented-implant loading data, (iii) improved preconditioning and scalable solvers for the multiphysics coupling, and (iv) extensions that leverage the smooth spectral machinery for gradient-based parameter estimation and model selection.
